\printedpage{29}
% head: A. Grothendieck : Dérivateurs, Ch. V, Édition des Universités de Montpellier II et Paris VII
% foot: Version 9 septembre 2001
$K_\infty$, whose union is $K_\infty$. But on the other hand I do not know how to give any convincing interpretation of $K_\infty$, like that of the ind-object of Cat\,hot defined by $\uuK = (K_n)_{n\in\mathbf{N}}$, which ind-represents the remarkable functor % [sic] defini -> défini, missing accent
\[
H = H_{\mathcal{M};x,y} : \Cathot^{o} \longrightarrow \Ens
\]
factoring the functor
\msp{page 45}
\[
X \longmapsto H_{\mathcal{M};x,y}(X) \defeq \Hom_{\mathcal{M}(X)\Sigma_X^{-1}}(x_X, y_X),
\]
by the formula, then
\[
H(X) = \varinjlim_{n} \underbrace{\pi_0\uHom(X, K_n)}_{\Hom_{\Cathot}(X,K_n)} .
\]
I do not know how to give an interpretation of this type either of the category $K_\infty$ itself, or of its image in Cat\,hot, or of its image in Hot. It is true that if I restrict the functor $H$ to the full subcategory of (Cat) consisting of \textit{finite} categories (Ob\,$X$ and Fl\,$X$ finite), then we will have
\[
\uHom^{!}(X, K_\infty) = \uHom(X, K_\infty),
\]
and the preceding isomorphism is interpreted as % [sic] précédant -> précédent, adjective intended
\[
\Hom_{\mathcal{M}(X)\Sigma_X^{-1}}(a_X, b_Y) \defeq H(X) \simeq \pi_0\uHom(X, K_\infty) \defeq \Hom_{\Cathot}(X, K_\infty).
\]
\texttt{[rather} $(a_X, b_X)$\texttt{]} This proves at least that knowledge of $K_\infty$ is useful. But this formula, with $X$ restricted to categories that are \textit{finite,} will not define $K_\infty$ as an object of Cat\,hot, up to unique isomorphism in this category - \textit{except} in the case where $K_\infty$ is itself homotopic to a
\msp{page 46}
finite category. (I do not quite see whether this happens very often \ldots\ But one cannot found a general theory of derivators and ``derivable'' categories of models on an assumption of this sort.)

But suppose that the functor $H$ under consideration even factors through Hot, i.e. that it transforms quasi-isomorphisms into bijections. We have seen that this is equivalent to the condition that it can be represented by a direct system over $\mathbf{N}$ made up of Thomason-Kan categories. It must moreover be easy, if one wants to, to arrange for the transition functors to be monomorphisms and fully faithful, so as to have a nice $\varinjlim$. It will of course be a Thomason-Kan category. Is this category determined by the functor $H : \Cathot^{o} \longrightarrow \Ens$ up to unique isomorphism in Cat\,hot? Of course, there is no question of its representing the functor $H$, except \texttt{[if]} the latter is really and truly representable, i.e. if the ind-object $\uuK$ of Cat\,hot is essentially constant (which must practically never be the case). Let us already note
\printedpage{30}
% head: A. Grothendieck : Dérivateurs, Ch. V, Édition des Universités de Montpellier II et Paris VII
% foot: Version 9 septembre 2001
\msp{page 47}
that if $u$ is a homomorphism of ind-objects \textit{in} Cat, or, what now amounts to the same thing, a homomorphism in Cat % [sic] homomorphimse -> homomorphisme, transposed letters
\[
K_\infty \xrightarrow{u_\infty} K'_\infty
\]
such that $\forall n$, $u_\infty(K_n) \subseteq K'_{n'}$ for a suitable $n' = \tau(n)$, and if we assume that the homomorphism of ind-objects induces an isomorphism for the ind-representable functors \textit{on} Cat\,hot (for example, if $\uuK$ and $\uuK'$ ind-represent the same functor $H$, and the homomorphism of ind-objects is the one corresponding to $\mathrm{id}_H$), then $u_\infty$ is necessarily an isomorphism \textit{in} Hot, hence also (unless I am mistaken) in Cat\,hot since $K_\infty$, $K'_\infty$ are Thomason-Kan categories. But what is not evident is the existence of a homomorphism $u_\infty$ as \texttt{[here-]}above, and if it exists, that it is unique (or at least canonical?) as an arrow of Cat\,hot.

For the existence of a $u_\infty$, we are reduced to proving that if we have a diagram with solid arrows
\[
\begin{tikzcd}
K_n \arrow[r, "i_n", "\text{cofib.}"'] \arrow[d, "u_n"] & K_{n+1} \arrow[dr, dashed] & \\
K'_{\tau(n)} \arrow[rr, "\text{transition}"'] & & K'_N ,
\end{tikzcd}
\]
and if an arrow $u_{n+1}$ is given ``in dotted lines'' \textit{in} Cat\,hot, so as to make the square commutative,
\msp{page 48}
then we can realize it by an arrow \textit{in} Cat that makes the diagram commute for real (in Cat). We are reduced to the same problem in a situation
\[
\begin{tikzcd}
X \arrow[r, hook, shift right, "i"'] & Y , \arrow[l, dashed, shift right, "p"']
\end{tikzcd}
\]
where $X$ is a Thomason category, $i$ a Thomason cofibration, $p$ a retraction from $Y$ to $X$ \textit{in} Cat\,hot, we want to realize it by a genuine retraction (\footnote[7]{\texttt{[a ?\ \ in the margin]}}).

Admitting that this is the case, hence that every homomorphism of the ind-objects $\uuK$, $\uuK'$ \textit{in} Cat\,hot can be represented by a homomorphism between the ind-objects \textit{in} Cat, hence also by a functor $K_\infty \longrightarrow K'_\infty$, there remains the question whether the corresponding homomorphism in Cat\,hot is unique. So if we have two functors $f$, $g$ which induce the same homomorphism on the ind-represented functors on Cat\,hot, whether $f$ and $g$ are homotopic. So

\textbf{Question}. Two functors $f, g : K_\infty \longrightarrow K'_\infty$ such that for every $n \in \mathbf{N}$, $f_n$, $g_n$ obtained from them by restriction to $K_n$ are homotopic (and even homotopic as functors from $K_n$ to a suitable $K'_N$), are they homotopic? (\footnote[8]{$K$, $K'$ Thomason categories, the $K_n$, $K'_n$ too.})

\printedpage{31}
% head: A. Grothendieck : Dérivateurs, Ch. V, Édition des Universités de Montpellier II et Paris VII
% foot: Version 9 septembre 2001
\msp{page 49}
I presume not - that it can happen that the minimum required length $l(n)$ for a path in $\uHom(K_n, K'_{\tau(n)})$ joining $f_n$ and $g_n$, tends to infinity if $n \longrightarrow \infty$. Does the hypothesis that the $K'_n$ are Thomason categories change anything here (are two homotopic objects in $\uHom(X, K')$, if $K'$ is Thomason, joined by an arrow?)

Of course, one can claim that the categories $\uCh_{\Sigma,n}(\mathcal{M}; a, b)$ have been defined ``canonically'', and their direct limit $\uCh_{\Sigma}(\mathcal{M}; a, b)$ too. But the definition was rather arbitrary and outlandish - I am convinced that ten mathematicians, proposing to develop this same idea of categories of paths, with the same aim as me, would arrive at ten different definitions! Yet it is not ruled out that we will end up finding a transitive system of isomorphisms in Cat\,hot, between the categories of ``infinite'' paths (i.e. of unspecified length) $K_\infty = K_\infty(\mathcal{M}; a, b)$ that we will have constructed, with the
\msp{page 50}
main aim (let us say) of representing the functor $H$ \textit{on finite categories.} And that the isomorphisms obtained will, of course, be homomorphisms in Ind(Cat\,hot). I have the impression that I have not yet thoroughly axiomatized the very notion of paths, which after all has a fairly obvious and important geometric meaning in its own right. So by what is to be understood by ``\textit{a} theory of paths in a category''. Such a theory $\mathcal{T}$ would thus associate to $\mathcal{M}$, $\Sigma$, $a$, $b$ a direct system of categories $\Ch^{\mathcal{T}}_{n,\Sigma}(\mathcal{M}; a, b)$. And the natural geometric constraints to be imposed on such theories, also with a view to their use in making explicit
\[
\Hom_{\mathcal{M}(X)\Sigma_X^{-1}}(a_X, b_X) \xrightarrow{\ \sim\ } \pi_0\uHom(X, \Chh^{\mathcal{T}}(\mathcal{M}; a, b)) % reading resolved: two underlines, the ind-path object.
\]
would give a transitive system of isomorphisms not only between the ind-objects of Cat\,hot, but perhaps also between the $\Ch^{\mathcal{T}}_{\infty,\Sigma}(\mathcal{M}; a, b)$, for different theories. If we have such a statement at our disposal, the category $\Ch_{\infty,\Sigma}(\mathcal{M}; a, b)$, or more exactly its images in Cat\,hot and Hot, % [sic] differentes -> différentes, missing accent
\msp{page 51}
would take on a convincing meaning.

On the other hand there is another point of view, which perhaps could ``canonify'' the categories $\uCh_\infty$ as an object of Cat\,hot and especially of Hot (one would pass to Cat\,hot by applying the functor $\psi$, Thomason-Kan envelope \ldots). Indeed, we are interested in constructing, in terms of $\mathcal{M}$, $\Sigma$, a ``derivator'' (or a prederivator) $\mathbf{D}_{\mathcal{M},\Sigma}$ especially on the ``domain'' of finite categories. Extending it to a larger domain is perhaps important in a ``second wind'', but can be considered a luxury in the first movement of founding a theory of derivators. This said, we can ask ourselves questions of the following type.

Let
\[
\Hot_0 \subset \Hot
\]
\printedpage{32}
% head: A. Grothendieck : Dérivateurs, Ch. V, Édition des Universités de Montpellier II et Paris VII
% foot: Version 9 septembre 2001
be the full subcategory of Hot, defined by finite categories (\footnote[9]{\textit{NB} Hot$_0$ is equivalent to a small category.}) - let us call the homotopy types ``of finite type'' (or of finite \textit{presentation}?). We have an obvious functor
\msp{page 52}
\[
\begin{aligned}
\Hot \quad &\longrightarrow \quad \Hot_0^{\wedge} = \uHom(\Hot_0^{o}, (\Ens)) \, ,\\
X \quad &\longmapsto \quad \tilde{X} = (\xi \longmapsto \Hom_{\Hot}(\xi, X)).
\end{aligned}
\]
This functor is \textit{conservative} \texttt{[this seems false]}, since a homomorphism $X \longrightarrow X'$ in Hot is an isomorphism if it induces isomorphisms for the functors $\tilde{X} \longrightarrow \tilde{X}'$, and even just for their values on the spheres $S^n$ ($n \geq 0$). (Unless I am mistaken. Are there any hoaxes, because the criterion for weak equivalence in terms of the $\pi_i$ does not apply directly, in the absence of base points. That should not be very serious \ldots\ \texttt{[this in fact seems serious]}). But is it perhaps even \textit{fully faithful} (\footnote[10]{\texttt{[a ?\ \ in the margin]}})? So that Hot would be realized as a full subcategory of $\Hot_0^{\wedge}$, containing the representable functors?